\subsubsection{Occlusion}
Occlusion, propuesta por Zeiler y Fergus, es una técnica de interpretabilidad que analiza la sensibilidad del modelo ocultando sucesivamente pequeñas regiones de la entrada y observando la variación de su salida. En imágenes, este análisis se realiza desplazando un cuadrado gris sobre la imagen; cuando el logit de la clase disminuye de forma notable, la región ``tapada'' se considera relevante para la decisión del clasificador.

En audio, adoptamos la misma idea pero aplicada sobre ventanas temporales o sobre regiones de un espectrograma, permitiendo identificar qué zonas de la señal contienen la información más importante a la hora de clasificar muestras.